% !TeX root = ../exemple-C-simulation.tex
% !TeX encoding = UTF-8
%% =====================================================================
%%  Exemple C — Compte rendu de simulation numérique (Python)
%%  Chaîne de travail :
%%     1.  python code/simulation.py      → écrit data/*.csv
%%     2.  compiler ce document           → tableaux, courbes et valeurs
%%                                          sont lus directement dans les CSV
%%  Montre : code Python inséré depuis le fichier réel (listings),
%%  tableau et courbes lus depuis des CSV (pgfplotstable / pgfplots),
%%  valeurs numériques du texte extraites des CSV, étude de convergence.
%% =====================================================================

% Valeurs lues dans data/resume.csv et data/tableau_periode.csv
\pgfplotstableread[col sep=comma]{data/resume.csv}\tabresume
\pgfplotstablegetelem{0}{ordre_observe}\of\tabresume \edef\OrdreObserve{\pgfplotsretval}
\pgfplotstablegetelem{0}{erreur_min}\of\tabresume \edef\ErreurMin{\pgfplotsretval}
\pgfplotstablegetelem{0}{erreur_max_pas_001}\of\tabresume \edef\ErreurMaxPas{\pgfplotsretval}
\pgfplotstableread[col sep=comma]{data/tableau_periode.csv}\tabperiode
\pgfplotstablegetelem{3}{T_rk4}\of\tabperiode \edef\TQuatreVingtDix{\pgfplotsretval}
\pgfplotstablegetelem{6}{T_rk4}\of\tabperiode \edef\TCentSoixanteDix{\pgfplotsretval}
\pgfplotstablegetelem{6}{T_serie}\of\tabperiode \edef\SerieCentSoixanteDix{\pgfplotsretval}

%% ---------------------------------------------------------------------
\begin{objectif}
%% ---------------------------------------------------------------------
Calculer numériquement la période $T$ d'un pendule simple lâché sans vitesse
d'un angle $\theta_{0}$ quelconque, comparer le résultat à la valeur exacte et
au développement en série, puis vérifier l'ordre de convergence de la méthode de
Runge–Kutta d'ordre~4 (RK4) utilisée.
\end{objectif}

%% ---------------------------------------------------------------------
\begin{modele}
%% ---------------------------------------------------------------------
L'angle $\theta$ d'un pendule simple de pulsation $\omega_{0}$ obéit à
$\ddot\theta+\omega_{0}^{2}\sin\theta=0$. Avec le temps réduit
$\tau=\omega_{0}t$, l'équation devient sans dimension :
\begin{equation}
  \dv[2]{\theta}{\tau} = -\sin\theta,
  \qquad \theta(0)=\theta_{0},\quad \dv{\theta}{\tau}(0)=0 .
  \label{eq:pendule}
\end{equation}
L'énergie réduite $\varepsilon=\tfrac12\bigl(\dv{\theta}{\tau}\bigr)^{2}-\cos\theta$
est conservée, ce qui fournit un test de la simulation. La période exacte
s'exprime à l'aide de l'intégrale elliptique complète $K$ ou de la
moyenne arithmético-géométrique $\mathrm{AGM}$~\cite{landau1976} :
\begin{equation}
  \frac{T}{T_{0}}=\frac{2}{\pi}\,K\bigl(\sin\tfrac{\theta_{0}}{2}\bigr)
  =\frac{1}{\mathrm{AGM}\bigl(1,\cos\tfrac{\theta_{0}}{2}\bigr)},
  \qquad T_{0}=2\pi .
  \label{eq:exacte}
\end{equation}
Pour les petites amplitudes, le développement en série s'écrit
\begin{equation}
  \frac{T}{T_{0}}=1+\frac{\theta_{0}^{2}}{16}+\frac{11\,\theta_{0}^{4}}{3072}
  +\bigO{\theta_{0}^{6}} .
  \label{eq:serie}
\end{equation}
\end{modele}

%% ---------------------------------------------------------------------
\begin{methode}
%% ---------------------------------------------------------------------
On pose $\vect{y}=(\theta,\ \mathrm{d}\theta/\mathrm{d}\tau)$ et
$\vect{f}(\vect{y})=(y_{2},\,-\sin y_{1})$. Un pas de RK4 de longueur $h$ est
\begin{equation}
  \vect{y}_{n+1}=\vect{y}_{n}+\frac{h}{6}\bigl(\vect{k}_{1}+2\vect{k}_{2}+2\vect{k}_{3}+\vect{k}_{4}\bigr),
  \label{eq:rk4}
\end{equation}
\begin{align}
  \vect{k}_{1}&=\vect{f}(\vect{y}_{n}),
  &\vect{k}_{2}&=\vect{f}\bigl(\vect{y}_{n}+\tfrac{h}{2}\vect{k}_{1}\bigr),\notag\\
  \vect{k}_{3}&=\vect{f}\bigl(\vect{y}_{n}+\tfrac{h}{2}\vect{k}_{2}\bigr),
  &\vect{k}_{4}&=\vect{f}\bigl(\vect{y}_{n}+h\,\vect{k}_{3}\bigr).
\end{align}
L'erreur globale de la méthode est en $\bigO{h^{4}}$. La période vaut quatre
fois la durée du premier passage par $\theta=0$, localisée par bissection sur le
dernier pas~(\cref{prog:rk4}).
\end{methode}

%% ---------------------------------------------------------------------
\begin{implementation}
%% ---------------------------------------------------------------------
Le \cref{prog:rk4} est extrait directement du fichier \code|code/simulation.py| :
toute modification du script est donc reprise à la prochaine compilation.
Le script complet (bibliothèque standard de Python uniquement) écrit les fichiers
\code|data/*.csv| lus ci-dessous.

\codefile[first-line=22, last-line=47]{Python}{code/simulation.py}%
  {Pas RK4 et calcul de la période (\texttt{code/simulation.py}).}{prog:rk4}
\end{implementation}

%% ---------------------------------------------------------------------
\begin{resultats}
%% ---------------------------------------------------------------------
\subsection{Période en fonction de l'amplitude}

Le \cref{tab:periode} compare la période calculée avec un pas $h=0{,}01$ à la
valeur exacte~\eqref{eq:exacte} et au développement~\eqref{eq:serie}. Pour
$\theta_{0}=\ang{90}$, on obtient $T/T_{0}=\num[round-mode=places,round-precision=5]{\TQuatreVingtDix}$.

\begin{table}[htbp]
  \caption{Période réduite $T/T_{0}$ (données : \texttt{data/tableau\_periode.csv}).}
  \label{tab:periode}
  \pgfplotstabletypeset[
    fixed, fixed zerofill,
    columns={theta0,T_rk4,T_exacte,T_serie,erreur_rel},
    columns/theta0/.style={column name={$\theta_{0}$ (\unit{\degree})}, int detect, column type=r},
    columns/T_rk4/.style={column name={RK4}, precision=6, column type=r},
    columns/T_exacte/.style={column name={Exacte}, precision=6, column type=r},
    columns/T_serie/.style={column name={Série~\eqref{eq:serie}}, precision=6, column type=r},
    columns/erreur_rel/.style={column name={Erreur relative RK4}, sci10, precision=1, column type=r},
  ]{data/tableau_periode.csv}
\end{table}

\begin{figure}[htbp]
  \begin{subfigure}[b]{0.49\textwidth}
    \begin{tikzpicture}
      \begin{axis}[width=0.95\linewidth, height=5.6cm, xlabel={$\theta_{0}$ (\unit{\degree})},
          ylabel={$T/T_{0}$}, xmin=0, xmax=180, ymin=0.9, ymax=3.4, grid=major,
          xtick={0,45,90,135,180}, ytick={1,2,3},
          legend pos=north west, scaled ticks=false]
        \addplot[thick, color=DataBlue] table[x=theta0, y=T_exacte, col sep=comma]{data/periode.csv};
        \addplot[only marks, mark=*, mark size=1.3pt, color=DataRed, mark repeat=2]
          table[x=theta0, y=T_rk4, col sep=comma]{data/periode.csv};
        \addplot[dashed, thick, color=TextColor] table[x=theta0, y=T_serie, col sep=comma]{data/periode.csv};
        \legend{Exacte, RK4, Série}
      \end{axis}
    \end{tikzpicture}
    \caption{Période réduite.}
    \label{fig:periode}
  \end{subfigure}\hfill
  \begin{subfigure}[b]{0.49\textwidth}
    \begin{tikzpicture}
      \begin{loglogaxis}[width=0.95\linewidth, height=5.6cm, xlabel={$h$},
          ylabel={Erreur relative sur $T$}, grid=major, legend pos=south east,
          xtick={0.01,0.1,1}, xticklabels={$10^{-2}$,$10^{-1}$,$1$},
          ytick={1e-10,1e-8,1e-6,1e-4}, yticklabels={$10^{-10}$,$10^{-8}$,$10^{-6}$,$10^{-4}$},
          minor tick num=0]
        \addplot[only marks, mark=*, mark size=1.5pt, color=DataRed]
          table[x=h, y=erreur_rel, col sep=comma]{data/convergence.csv};
        \addplot[dashed, thick, color=TextColor, domain=0.0125:0.4, samples=2]
          {1.345e-10*(x/0.0125)^4};
        \legend{RK4, pente $4$}
      \end{loglogaxis}
    \end{tikzpicture}
    \caption{Convergence en $h$ ($\theta_{0}=\ang{60}$).}
    \label{fig:convergence}
  \end{subfigure}
  \caption{Résultats de la simulation, tracés directement depuis les fichiers
    CSV : (a) période en fonction de l'amplitude ; (b) erreur relative de RK4
    en fonction du pas, en échelle logarithmique.}
  \label{fig:resultats}
\end{figure}
\end{resultats}

%% ---------------------------------------------------------------------
\begin{validation}
%% ---------------------------------------------------------------------
Pour $\theta_{0}=\ang{60}$, l'erreur relative décroît de
$\num{1.3e-4}$ ($h=0{,}4$) à $\num[round-mode=figures,round-precision=2]{\ErreurMin}$
($h=0{,}0125$). La pente de $\log(\text{erreur})$ en fonction de $\log h$, ajustée
sur les six pas, vaut
\begin{equation}
  p=\num{\OrdreObserve},
  \label{eq:ordre}
\end{equation}
en accord avec l'ordre théorique $4$ de RK4 (\cref{fig:convergence}).
\end{validation}

%% ---------------------------------------------------------------------
\begin{analyse}
%% ---------------------------------------------------------------------
\paragraph{Qualité de la simulation.}
Avec $h=0{,}01$, l'erreur relative sur $T$ reste inférieure à
$\num[round-mode=figures,round-precision=1]{\ErreurMaxPas}$ pour toutes les amplitudes
de \ang{5} à \ang{170} : la précision est limitée par l'arithmétique en double
précision plutôt que par la méthode.

\paragraph{Validité du développement en série.}
La série~\eqref{eq:serie} reproduit la période à mieux que \num{1e-3} jusqu'à environ
\ang{60}, mais s'en écarte fortement ensuite : pour $\theta_{0}=\ang{170}$ elle
donne $T/T_{0}=\num[round-mode=places,round-precision=2]{\SerieCentSoixanteDix}$ contre
$\num[round-mode=places,round-precision=2]{\TCentSoixanteDix}$ pour la valeur exacte, soit une
sous-estimation d'environ \qty{25}{\percent}. La période diverge lorsque
$\theta_{0}\to\pi$ (pendule lâché au voisinage de la position d'équilibre
instable), ce qu'aucun développement polynomial ne peut représenter.
\end{analyse}

%% ---------------------------------------------------------------------
\begin{conclusion}
%% ---------------------------------------------------------------------
L'intégration RK4 reproduit la période exacte du pendule non linéaire avec un
ordre de convergence mesuré $p=\num{\OrdreObserve}$. Le développement en série
n'est fiable que pour les amplitudes modérées. La chaîne
« script Python $\to$ fichiers CSV $\to$ document \LaTeX{} » garantit que les
tableaux, les courbes et les valeurs citées proviennent toujours de la
dernière exécution du programme.
\end{conclusion}
